
\chapter{Yêu cầu dữ liệu}

Chương này mô tả các yêu cầu về dữ liệu của hệ thống, bao gồm các nhóm dữ liệu chính, mối quan hệ giữa các thực thể và các yêu cầu về tính toàn vẹn, lưu trữ và bảo mật dữ liệu.

%%%%%%%%%%%%%%%%%%%%%%%%%%%%%%%%%%%%%%%%%%%%%%%%%%
\section{Tổng quan dữ liệu}

Hệ thống quản lý dịch vụ bảo trì máy lọc nước sử dụng cơ sở dữ liệu quan hệ (Relational Database) nhằm lưu trữ toàn bộ thông tin phục vụ các nghiệp vụ quản lý.

Dữ liệu được tổ chức theo nhiều thực thể (Entity) có mối quan hệ với nhau, đảm bảo tính nhất quán, khả năng mở rộng và hỗ trợ truy xuất dữ liệu hiệu quả.

Các nhóm dữ liệu chính bao gồm:

\begin{itemize}
    \item Thông tin người dùng và phân quyền.
    \item Thông tin sản phẩm.
    \item Thông tin tồn kho.
    \item Thông tin lô hàng.
    \item Thông tin khách hàng.
    \item Thông tin hợp đồng.
    \item Thông tin thiết bị.
    \item Dữ liệu Dashboard thiết bị.
    \item Thông tin nhân viên.
    \item Nhật ký hoạt động (Activity Log).
\end{itemize}

%%%%%%%%%%%%%%%%%%%%%%%%%%%%%%%%%%%%%%%%%%%%%%%%%%
\section{Các thực thể dữ liệu}

Bảng \ref{tab:data_entities} mô tả các thực thể chính của hệ thống.

\begin{table}[H]
\centering
\caption{Danh sách các thực thể dữ liệu}
\label{tab:data_entities}

\begin{tabularx}{\textwidth}{|l|X|}
\hline
\textbf{Thực thể} & \textbf{Mô tả} \\
\hline

User &
Lưu thông tin tài khoản và phân quyền người dùng. \\
\hline

Employee &
Lưu thông tin nhân viên và kỹ thuật viên. \\
\hline

Customer &
Lưu thông tin khách hàng sử dụng dịch vụ. \\
\hline

Product &
Lưu thông tin các dòng máy lọc nước và linh kiện. \\
\hline

Inventory &
Quản lý số lượng sản phẩm hiện có trong kho. \\
\hline

Batch &
Quản lý các lô hàng nhập kho. \\
\hline

Contract &
Lưu thông tin hợp đồng lắp đặt, bảo trì và thay thế thiết bị. \\
\hline

Device &
Lưu thông tin các thiết bị đã được lắp đặt cho khách hàng. \\
\hline

Device Dashboard &
Lưu dữ liệu thống kê và dữ liệu cảm biến của thiết bị. \\
\hline

Activity Log &
Ghi nhận lịch sử thao tác của người dùng trên hệ thống. \\
\hline

\end{tabularx}

\end{table}

%%%%%%%%%%%%%%%%%%%%%%%%%%%%%%%%%%%%%%%%%%%%%%%%%%
\section{Mối quan hệ dữ liệu}

Các thực thể trong hệ thống được liên kết với nhau nhằm đảm bảo dữ liệu được quản lý thống nhất.

Một số mối quan hệ chính bao gồm:

\begin{itemize}
    \item Một khách hàng có thể sở hữu nhiều thiết bị.
    \item Một khách hàng có thể có nhiều hợp đồng.
    \item Một hợp đồng được gắn với một khách hàng.
    \item Một thiết bị thuộc về một khách hàng.
    \item Một sản phẩm có thể xuất hiện trong nhiều lô hàng.
    \item Một sản phẩm được quản lý trong tồn kho.
    \item Một nhân viên có thể thực hiện nhiều hoạt động bảo trì.
    \item Một người dùng có thể tạo nhiều bản ghi Activity Log.
\end{itemize}

Hình \ref{fig:erd} minh họa sơ đồ thực thể - liên kết (Entity Relationship Diagram) của hệ thống.

\begin{figure}[H]
    \centering
    \includegraphics[width=\textwidth]{figures/database/erd.png}
    \caption{Sơ đồ ERD của hệ thống}
    \label{fig:erd}
\end{figure}

%%%%%%%%%%%%%%%%%%%%%%%%%%%%%%%%%%%%%%%%%%%%%%%%%%
\section{Yêu cầu lưu trữ dữ liệu}

Hệ thống sử dụng MySQL làm hệ quản trị cơ sở dữ liệu.

Các yêu cầu lưu trữ bao gồm:

\begin{itemize}
    \item Dữ liệu được lưu trữ tập trung trong cơ sở dữ liệu MySQL.
    \item Mỗi bản ghi phải có khóa chính (Primary Key) để định danh duy nhất.
    \item Các bảng có quan hệ phải sử dụng khóa ngoại (Foreign Key) nhằm đảm bảo tính toàn vẹn tham chiếu.
    \item Dữ liệu phải được chuẩn hóa nhằm hạn chế dư thừa và tránh bất nhất dữ liệu.
    \item Hệ thống phải hỗ trợ sao lưu (Backup) và khôi phục (Restore) dữ liệu khi cần thiết.
\end{itemize}

%%%%%%%%%%%%%%%%%%%%%%%%%%%%%%%%%%%%%%%%%%%%%%%%%%
\section{Yêu cầu về tính toàn vẹn dữ liệu}

Để đảm bảo tính chính xác và nhất quán của dữ liệu, hệ thống cần đáp ứng các yêu cầu sau:

\begin{itemize}
    \item Không cho phép trùng lặp mã định danh của các thực thể.
    \item Không cho phép tạo hợp đồng khi chưa tồn tại khách hàng.
    \item Không cho phép tạo thiết bị khi chưa có sản phẩm tương ứng.
    \item Dữ liệu nhập vào phải được kiểm tra tính hợp lệ trước khi lưu.
    \item Các thao tác thêm, sửa và xóa dữ liệu phải được ghi nhận trong Activity Log.
\end{itemize}

%%%%%%%%%%%%%%%%%%%%%%%%%%%%%%%%%%%%%%%%%%%%%%%%%%
\section{Yêu cầu bảo mật dữ liệu}

Dữ liệu của hệ thống cần được bảo vệ nhằm đảm bảo tính bảo mật và an toàn thông tin.

Các yêu cầu bao gồm:

\begin{itemize}
    \item Chỉ người dùng đã xác thực mới được phép truy cập hệ thống.
    \item Dữ liệu được phân quyền theo từng nhóm người dùng.
    \item Mật khẩu người dùng phải được mã hóa trước khi lưu trong cơ sở dữ liệu.
    \item Chỉ Administrator mới có quyền quản lý tài khoản và phân quyền.
    \item Hệ thống phải ghi nhận các thao tác quan trọng thông qua Activity Log để phục vụ việc kiểm tra và truy vết.
\end{itemize}

%%%%%%%%%%%%%%%%%%%%%%%%%%%%%%%%%%%%%%%%%%%%%%%%%%
\section{Chi tiết 17 Database Tables}

Hệ thống hiện tại được xây dựng với 17 bảng dữ liệu trong MySQL, được tạo thông qua Eloquent Migrations của Laravel. Mỗi bảng được thiết kế để lưu trữ một loại thực thể cụ thể với các khóa ngoại để đảm bảo tính toàn vẹn dữ liệu.

\begin{table}[H]
\centering
\caption{17 Database Tables chi tiết}
\label{tab:database_tables}
\small
\begin{tabularx}{\textwidth}{|l|l|c|X|}
\hline
\textbf{Bảng} & \textbf{Model} & \textbf{Records} & \textbf{Mô tả} \\
\hline

roles & Role & 3 & Administrator, Employee, Technician \\
\hline

users & User & 3 & Tài khoản người dùng (bcrypt password) \\
\hline

employees & Employee & 9 & Thông tin nhân viên với FK to roles \\
\hline

categories & Category & 4 & Danh mục sản phẩm (RO, Nano, Industrial) \\
\hline

products & Product & 10 & Sản phẩm máy lọc nước (AQ-RO-50 to AQ-CORE-RO) \\
\hline

suppliers & Supplier & 3 & Nhà cung cấp hàng \\
\hline

inventories & Inventory & 10 & Tồn kho (1:1 với products) \\
\hline

batches & Batch & 6 & Lô hàng nhập (LOT-2025-001 to LOT-2025-006) \\
\hline

batch\_details & BatchDetail & 30+ & Chi tiết lô hàng (many-to-many) \\
\hline

customers & Customer & 24 & Khách hàng (soft deletes enabled) \\
\hline

contracts & Contract & 18 & Hợp đồng bảo trì \\
\hline

contract\_services & ContractService & 3 & Dịch vụ tham chiếu (static) \\
\hline

devices & Device & 30 & Thiết bị (TB-01000 to TB-01029) \\
\hline

device\_dashboard\_data & DeviceDashboardData & 24+ & Dữ liệu cảm biến IoT (TDS, temp, flow, pH) \\
\hline

maintenance\_records & MaintenanceRecord & 20 & Phiếu bảo trì (BT-0001 to BT-0020) \\
\hline

activity\_logs & ActivityLog & 25+ & Nhật ký hoạt động người dùng \\
\hline

attachments & Attachment & 0 & Tệp đính kèm (polymorphic) \\
\hline

\end{tabularx}

\end{table}

\subsection*{Tổng cộng dữ liệu được seeded}

Hệ thống hiện tại được populate với hơn 100 records dữ liệu test:

\begin{itemize}
    \item 3 Roles + 9 Employees + 3 Users
    \item 4 Categories + 10 Products + 10 Inventories
    \item 3 Suppliers + 6 Batches + 30 Batch Details
    \item 24 Customers + 18 Contracts + 3 Contract Services
    \item 30 Devices + 24 Device Telemetry Data
    \item 20 Maintenance Records + 25 Activity Logs
\end{itemize}

%%%%%%%%%%%%%%%%%%%%%%%%%%%%%%%%%%%%%%%%%%%%%%%%%%
\section{REST API Endpoints}

Bên cạnh giao diện Web, hệ thống cung cấp 30+ REST API endpoints cho phép các ứng dụng bên ngoài truy cập dữ liệu:

\begin{itemize}
    \item \textbf{Products API}: GET, POST, PUT, DELETE /api/products
    \item \textbf{Categories API}: GET, POST, PUT, DELETE /api/categories
    \item \textbf{Inventories API}: GET, PATCH /api/inventories
    \item \textbf{Customers API}: GET, POST, PUT, DELETE /api/customers
    \item \textbf{Devices API}: GET, POST /api/devices
    \item \textbf{Contracts API}: GET, POST /api/contracts
\end{itemize}

Mỗi endpoint trả về JSON responses với proper HTTP status codes (200, 201, 422, 404, 500).

